\lysh{34816}{4}
\begin{alist}
\item Η γωνία του κυκλικού τομέα που αντιστοιχεί στους εργαζομένους της κατηγορίας Β$'$ γνωρίζουμε ότι είναι $135^\circ$. Το κυκλικό διάγραμμα αντιστοιχεί σε $360^\circ$, συνεπώς η κατηγορία Β$'$ θα είναι το $\frac{135}{360} = 0,375$ του συνόλου ή από τον τύπο $a_i = \nu_i \cdot 360^\circ$ και $f_i$ για την κατηγορία Β$'$ έχουμε: $135^\circ = 360^\circ \cdot f_B \Leftrightarrow f_B = \frac{135^\circ}{360^\circ} = 0,375$.
Συνεπώς οι εργαζόμενοι που ανήκουν στην κατηγορία Β$'$ θα είναι το $37,5\%$ του συνόλου δηλαδή $0,375 \cdot 1200 = 450$ εργαζόμενοι.

\item Οι εργαζόμενοι που ανήκουν στην κατηγορία Α$'$ αντιστοιχούν στο $25\%$ του συνόλου δηλαδή $0,25 \cdot 1200 = 300$ εργαζόμενοι.
Οι εργαζόμενοι που ανήκουν στην κατηγορία Δ$'$ αντιστοιχούν στο $12,5\%$ του συνόλου δηλαδή οι μισοί της κατηγορίας Α$'$ άρα $150$ εργαζόμενοι.
Όλοι μαζί είναι $1200$ και μέχρι τώρα έχουμε βρει:
\begin{itemize}
	\item Κατηγορία Α$'$ = 300
	\item Κατηγορία Β$'$ = 450
	\item Κατηγορία Γ$'$ = ....
	\item Κατηγορία Δ$'$ = 150
	\item Κατηγορία Ε$'$ = ....
\end{itemize}
Γνωρίζουμε όμως ότι οι εργαζόμενοι της κατηγορίας Γ$'$ (έστω $2x$) είναι διπλάσιοι από αυτούς της κατηγορίας Ε$'$ (έστω $x$) συνεπώς θα έχουμε:
\[ 300+450+2x+150+x=1200 \Leftrightarrow 900+3x=1200 \Leftrightarrow 3x=1200-900 \Leftrightarrow 3x=300 \Leftrightarrow x=100, \]
άρα:
\begin{itemize}
	\item Κατηγορία Γ$'$ = 200
	\item Κατηγορία Ε$'$ = 100
\end{itemize}

\item 
\begin{rlist}
	\item Από το β) έχουμε για το πλήθος των εργαζομένων σε κάθε κατηγορία τα εξής:
	\begin{itemize}
		\item Κατηγορία Α$'$ $300$ εργαζόμενοι, άρα $\nu_1=300$.
		\item Κατηγορία Β$'$ $450$ εργαζόμενοι, άρα $\nu_2=450$.
		\item Κατηγορία Γ$'$ $200$ εργαζόμενοι, άρα $\nu_3=200$.
		\item Κατηγορία Δ$'$ $150$ εργαζόμενοι, άρα $\nu_4=150$.
		\item Κατηγορία Ε$'$ $100$ εργαζόμενοι, άρα $\nu_5=100$.
	\end{itemize}
	Με βάση τα παραπάνω προκύπτει το ραβδόγραμμα συχνοτήτων $\nu_i$.
	\begin{center}
%	\begin{tikzpicture}
%	\begin{axis}[
%	    axis lines = middle,
%	    ylabel = {$\nu_i$},
%	    ymin = 0, ymax = 550,
%	    ytick = {0,50,100,150,200,250,300,350,400,450,500,550},
%	    grid = both,
%	    grid style = {gray!30},
%	    width = 12cm, height = 8cm,
%	    tick label style = {font=\scriptsize},
%	    every axis y label/.style = {at={(current axis.above origin)}, anchor=south},
%	    ybar,
%	    bar width=20pt,
%	    enlarge x limits=0.15,
%	    xtick=data,
%	    xticklabels={Κατηγορία Α$'$, Κατηγορία Β$'$, Κατηγορία Γ$'$, Κατηγορία Δ$'$, Κατηγορία Ε$'$},
%	    x tick label style={align=center},
%	]
%	    \addplot[very thick, maincolor, fill=maincolor!30] coordinates {
%	        (Κατηγορία Α$'$, 300)
%	        (Κατηγορία Β$'$, 450)
%	        (Κατηγορία Γ$'$, 200)
%	        (Κατηγορία Δ$'$, 150)
%	        (Κατηγορία Ε$'$, 100)
%	    };
%	\end{axis}
%	\end{tikzpicture}
	\end{center}
	
	\item Όταν συνταξιοδοτηθούν οι εργαζόμενοι της κατηγορίας Γ΄ που στο σύνολό τους είναι $200$, τότε στην εταιρεία θα απασχολούνται $1200-200=1000$ εργαζόμενοι. Οι εργαζόμενοι της κατηγορίας Ε΄ γνωρίζουμε από το β) ότι είναι $100$ συνεπώς αντιπροσωπεύουν το $\frac{100}{1000} = 0,1$ του συνόλου των εργαζομένων, δηλαδή το $10\%$.
\end{rlist}
\end{alist}
\vspace{6mm}
\noindent\rule{\textwidth}{0.4pt}
\vspace{6mm}

